% !TEX program = xelatex
\documentclass[12pt,a4paper]{ctexart}

% ================================================================
%  SCX Prize Declaration:
%  ——
%  The SCX Prize Declaration:
%  An Auditable Alternative to the Nobel Prize — Redefining
%  Scientific Excellence Under Declared Audit Standards
%  Chinese + English bilingual  800+ lines  Calm, mathematical, final
% ================================================================

% Packages
\usepackage{amsmath,amssymb,amsthm}
\usepackage{mathtools}
\usepackage{bm}
\usepackage{graphicx}
\usepackage{hyperref}
\usepackage{geometry}
\usepackage{enumitem}
\usepackage{booktabs}
\usepackage{tikz}
\usepackage{xcolor}
\usepackage{cleveref}
\usepackage{bbm}
\usepackage{float}
\usepackage{physics}
\usepackage{setspace}
\usepackage{longtable}
\usepackage{multirow}
\usepackage{array}
\usepackage{makecell}
\usepackage{dsfont}

\geometry{margin=2.5cm}
\onehalfspacing

% Hyperref setup
\hypersetup{
    colorlinks=true,
    linkcolor=blue,
    citecolor=blue,
    urlcolor=blue,
    pdftitle={SCX Prize Declaration——},
    pdfauthor={SCX},
    pdfsubject={SCX Prize Declaration},
}

% ================================================================
% Theorem environments
% ================================================================
\newtheorem{theorem}{Theorem}[section]
\newtheorem{lemma}[theorem]{Lemma}
\newtheorem{proposition}[theorem]{Proposition}
\newtheorem{corollary}[theorem]{Corollary}
\newtheorem{definition}[theorem]{Definition}
\newtheorem{conjecture}[theorem]{Conjecture}
\newtheorem{axiom}[theorem]{Axiom}

\theoremstyle{remark}
\newtheorem{remark}[theorem]{Remark}
\newtheorem{example}[theorem]{Example}

% ================================================================
% Custom math commands
% ================================================================
\newcommand{\R}{\mathbb{R}}
\newcommand{\C}{\mathbb{C}}
\newcommand{\N}{\mathbb{N}}
\newcommand{\Z}{\mathbb{Z}}
\newcommand{\E}{\mathbb{E}}
\newcommand{\Pbb}{\mathbb{P}}
\newcommand{\Var}{\mathrm{Var}}
\newcommand{\Cov}{\mathrm{Cov}}
\newcommand{\indep}{\perp\!\!\!\perp}

% SCX-specific operators
\newcommand{\SCX}{\	extsf{SCX}}
\newcommand{\SCXPrize}{\mathsf{SCX\text{-}Prize}}
\newcommand{\gaugeVec}{\mathbf{g}}
\newcommand{\sumgd}{\sum_m \mathbf{g}_m = \mathbf{0}}
\newcommand{\gaugeGroup}{\mathbb{G}}
\newcommand{\auditOp}{\mathcal{A}}
\newcommand{\scoreOp}{\mathcal{S}}
\newcommand{\prizeSet}{\mathcal{P}}
\newcommand{\candidateSet}{\mathcal{C}}
\newcommand{\criterionSet}{\mathcal{K}}

\DeclareMathOperator{\argmax}{arg\,max}
\DeclareMathOperator{\argmin}{arg\,min}
\DeclareMathOperator{\sgn}{sgn}
\DeclareMathOperator{\diag}{diag}
\DeclareMathOperator{\spec}{spec}
\DeclareMathOperator{\rank}{rank}
\DeclareMathOperator{\Audit}{Audit}
\DeclareMathOperator{\Score}{Score}
\DeclareMathOperator{\Prize}{Prize}
\DeclareMathOperator{\Declare}{Declare}

% ================================================================
% Honest critique markers
% ================================================================
\newcommand{\honestcrit}[1]{%
    \textcolor{red}{\textbf{#1}}%
}

\newenvironment{honestbox}{%
    \medskip\noindent
    \color{red}
    \begin{minipage}{0.92\textwidth}
    \textbf{}
    \begin{itemize}[leftmargin=*]
}{%
    \end{itemize}
    \end{minipage}
    \medskip
}

% Review/audit note
\newcommand{\auditnote}[1]{%
    \medskip\noindent
    \textbf{ (Audit Note):} \textit{#1}%
}

% Proof-status markers
\newcommand{\rigorFull}{\textnormal{\textbf{[]}}}
\newcommand{\rigorPartial}{\textnormal{\textbf{[]}}}
\newcommand{\rigorSketch}{\textnormal{\textbf{[]}}}
\newcommand{\openproblem}{\textnormal{\textbf{[]}}}

% ================================================================
% Title
% ================================================================
\title{
    \textbf{\LARGE SCX\\[6pt]
    ——}\\[14pt]
    \large
    The SCX Prize Declaration\\[4pt]
    An Auditable Alternative to the Nobel Prize:\\[2pt]
    Redefining Scientific Excellence Under Declared Audit Standards\\[10pt]
    \normalsize
    ---  ---\\[2pt]
    --- This is not a petition. It is a theorem's consequence. ---
}

\author{%
    SCX \\[2pt]
    \textit{The SCX Audit Council}\\[6pt]
    \small \\[2pt]
    \small Gauge-Fixing the Comparative Legitimacy of Scientific Excellence
}

\date{202671 \\[2pt] July 1, 2026}

\begin{document}

\maketitle

% ================================================================
% ABSTRACT
% ================================================================
\begin{abstract}

\noindent\textbf{.}
——$M=1$SCXSCX——$g$SCX$\sumgd$API——SCX

\medskip
\noindent\textbf{English Abstract.}
Any institution claiming to award excellence must first declare its audit standards. Institutions that do not declare standards make claims that are mathematically void — they perform $M=1$ evaluations against an undefined function and announce the result as universal truth. This paper establishes the SCX Prize: a fully auditable scientific excellence prize system in which winners are determined by a public, reproducible algorithm with no committee. The SCX Prize has eight categories, each with a computable audit criterion. The prize's legitimacy derives not from the authority of a deliberation committee but from the publicity and reproducibility of its standards. Proceeding from gauge-fixing theory, we prove that Nobel Prize committees are gauge-unequivalent to the fields they judge — their $g$ parameters are undeclared, rendering their claims mathematically void. The SCX Prize, by defining the fixed gauge condition $\sumgd$, establishes the comparative legitimacy of excellence awards for the first time. We also declare the prize's funding source (a fixed percentage of Yajie API revenue) and its self-audit mechanism — the SCX Prize submits to the same audit standard it imposes on others.

\end{abstract}

\newpage

% ================================================================
% TABLE OF CONTENTS
% ================================================================
\tableofcontents
\newpage

% ================================================================
% PREAMBLE: THEOREM ZERO
% ================================================================
\section*{}
\label{sec:preamble}
\addcontentsline{toc}{section}{ (Preamble: Theorem Zero)}

\begin{quotation}
\noindent\textit{``——SCX——"}

\medskip
\noindent\textit{``If a claim cannot be audited, it is not even a claim. It is noise — a sound wave that transmits energy through a communication channel but carries zero information. The Nobel Prize claims it awards excellence. The SCX Prize does not claim — it declares standards, then executes a computation. The difference is that the former substitutes authority for audit; the latter substitutes audit for authority.''}
\end{quotation}

\medskip
\noindent
\auditnote{This document is written in the calm, declarative tone of a mathematical result. It does not petition. It does not persuade. It states the consequence of a theorem and then describes the institutional mechanism that follows from it. The reader is free to disagree with the theorem's premises — but the conclusion follows from them with the inexorability of modus ponens.}

\bigskip

\begin{center}
\fbox{%
\begin{minipage}{0.88\textwidth}
\centering
\textbf{ 0 ()}\\[4pt]
\textbf{Axiom 0 (The Auditability of Excellence)}

\vspace{6pt}
 $\mathcal{E}$  $\auditOp: \mathcal{E} \to \{\text{}, \text{}, \bot\}$  $\bot$ ``''

\medskip
\textit{Chinese:}  $\auditOp(e) = \bot$ $e$ ——

\medskip
\textit{English:} If $\auditOp(e) = \bot$, then $e$ is not an excellence claim — it is an excellence assertion. The difference between a claim and an assertion is: claims accept audit; assertions refuse audit. Only claims have mathematical meaning.

\vspace{4pt}
$\boxed{\auditOp(e) \neq \bot \iff e \in \text{}}$
\end{minipage}%
}
\end{center}

\newpage

% ================================================================
% §1. THE AUDIT VACUUM IN PRIZE-GIVING
% ================================================================
\section{$M=1$}
\label{sec:audit_vacuum}

\subsection{}

\noindent
\textbf{1.1  (Prizes as Claims)}

————\textbf{} $\candidateSet$  $w \in \candidateSet$ 

\begin{equation}
    w = \argmax_{c \in \candidateSet} \; \text{}(c),
    \label{eq:prize_claim}
\end{equation}

$(\cdot)$——\textbf{$(\cdot)$}

$(\cdot)$deliberation\textit{}

——\textbf{}

\begin{honestbox}
    \item 5050————\textit{}
    \item 
    \item ACM``''——``''
\end{honestbox}

\subsection{$M=1$ }

\noindent
\textbf{1.2 $M=1$  (Formalizing $M=1$ Undeclared-ness)}

\begin{definition}[, Prize Function]
\label{def:prize_function}
 $\prizeSet$ $\prizeSet$ 
\begin{equation}
    \prizeSet = (\candidateSet, \scoreOp, \gaugeVec, \text{procedure}, \text{declaration}),
    \label{eq:prize_quintuple}
\end{equation}

\begin{itemize}
    \item $\candidateSet$
    \item $\scoreOp: \candidateSet \to \R$
    \item $\gaugeVec = (\mathbf{g}_1, \dots, \mathbf{g}_M)$
    \item $\text{procedure}$
    \item $\text{declaration}$
\end{itemize}
\end{definition}

\begin{definition}[$M=1$ , $M=1$ Undeclared Prize]
\label{def:M1_undeclared}
 $\prizeSet$  \textbf{$M=1$ }
\begin{enumerate}[label=(\roman*)]
    \item $\gaugeVec$ 
    \item $\scoreOp$  $\candidateSet$  $\R$ 
    \item $\text{declaration} = \text{false}$
\end{enumerate}
\end{definition}

\noindent
\auditnote{The term ``$M=1$'' is not a typo. A committee of five people whose $g$ parameters are undeclared is mathematically equivalent to a single evaluator whose $g$ is unknown. The ``committee'' is not an audit structure — it is a black-box oracle whose internal state is inaccessible. From an information-theoretic standpoint, $M$ undeclared evaluators $\equiv$ 1 undeclared evaluator. The $M$ is not a multiplicity of perspectives — it is a multiplicity of hidden parameters.}

\begin{theorem}[$M=1$ , \\
    Information-Theoretic Voidness of $M=1$ Undeclared Prizes]
\label{thm:M1_void}
 $\prizeSet$  $M=1$  $w$``$w$ '' $I(w; \text{truth})$ 
\begin{equation}
    I(w; \text{truth} \mid \gaugeVec \text{ undeclared}) \leq 0.
    \label{eq:info_zero}
\end{equation}
\end{theorem}

\begin{proof}[ (Proof Sketch)]
 $I(w; \text{truth})$ ``'' $\gaugeVec$  $\scoreOp_g(c)$  $c$  $\gaugeVec$  $\gaugeVec$  $\gaugeVec$ $\scoreOp_g$ ``''—— $\gaugeVec$ $\scoreOp_g(c)$  $\text{truth}(c)$  $\gaugeVec$ 
\end{proof}

\begin{corollary}[]
\label{cor:nobel_info}
``''\textit{}——\textit{}
\end{corollary}

\subsection{}

\noindent
\textbf{1.3  (Committees Are Not Audit Structures)}

``''$M$

\begin{enumerate}[label=(\alph*)]
    \item \textbf{}\textit{}\textit{}
    \item \textbf{}——
    \item \textbf{}
\end{enumerate}

\auditnote{A committee of five experts is not five independent audits. It is one conversation with five participants. The output is a consensus — the annihilation of minority viewpoints. The SCX audit standard requires that all $M$ evaluators' scores be public and independently verifiable. Consensus is not an audit product; it is an audit failure — it destroys the diversity that makes $M > 1$ meaningful.}

\begin{honestbox}
    \item \textbf{}qualitative deliberation
    \item $g$self-referential system0
    \item ——\textit{}
\end{honestbox}

\subsection{}

\noindent
\textbf{1.4  (Historical Note: Arrogance as Gauge Freedom)}

1901————

SCX\textbf{}$g$——

\auditnote{The Nobel Foundation's 50-year secrecy rule is not an administrative detail. It is a gauge-theoretic structure: it protects the committee's $g$ parameter from ever being audited. By the time the records are opened, all participants are dead, the field has moved on, and the audit is a historical curiosity rather than a live accountability mechanism. This is not an accident. It is a design.}

\newpage

% ================================================================
% §2. THE SCX PRIZE MECHANISM
% ================================================================
\section{SCX}
\label{sec:mechanism}

\subsection{}

\noindent
\textbf{2.1  (Core Idea: No Human Judges)}

SCX

\begin{center}
\fbox{%
\begin{minipage}{0.88\textwidth}
\centering
\textbf{SCX (SCX Prize First Principle)}

\vspace{4pt}
50

\medskip
\textit{No committee. No human judgment. No deliberation. No voting. No negotiation. No consensus. No majority opinion. No minority opinion. No sealed archives. No 50-year declassification period.}

\vspace{4pt}


\medskip
\textit{The algorithm runs. The highest score wins. Anyone can re-run the algorithm and verify the result.}
\end{minipage}%
}
\end{center}

\subsection{}

\noindent
\textbf{2.2  (Prize Architecture)}

SCX

\begin{enumerate}[label=\textbf{P\arabic*}.]
    \item \textbf{} $\candidateSet$——
    
    \item \textbf{} $\scoreOp_k: \candidateSet \to \R$ $k \in \{1, \dots, 8\}$
    
    \item \textbf{} $\sumgd$ $\gaugeVec$ ``''
    
    \item \textbf{}
    
    \item \textbf{}
\end{enumerate}

\subsection{}

\noindent
\textbf{2.3  (Formalizing Algorithmic Judgment)}

\begin{definition}[, Auditable Prize]
\label{def:auditable_prize}
\textbf{} $\prizeSet_{\text{auditable}}$ 
\begin{enumerate}[label=(\roman*)]
    \item $\scoreOp$ 
    \item $\gaugeVec$ 
    \item $\text{declaration} = \text{true}$
    \item  $c \in \candidateSet$$\scoreOp(c)$ 
\end{enumerate}
\end{definition}

\begin{proposition}[]
\label{prop:auditable_info}
 $\prizeSet_{\text{auditable}}$ ``$w$  $k$ ''
\begin{equation}
    I(w; \text{truth} \mid \gaugeVec \text{ declared}) > 0,
    \label{eq:info_positive}
\end{equation}
 $\scoreOp$ 
\end{proposition}

\begin{proof}[ (Proof Sketch)]
 $\gaugeVec$  $\scoreOp(c)$  $c$  $\scoreOp$ $w = \argmax_c \scoreOp(c)$  $\scoreOp$ $w$  $\scoreOp$ —— $\scoreOp$ 
\end{proof}

\auditnote{This is the crucial distinction. The SCX Prize does not claim to identify the ``truly best'' scientist in some metaphysical sense. It claims: ``Under the publicly declared scoring function $\scoreOp_k$, candidate $w$ achieved the highest score.'' The truth of this claim is trivially verifiable — just re-run the computation. The debate then shifts to whether $\scoreOp_k$ is a reasonable measure of excellence — and \textit{that} debate is healthy, public, and improvable. The Nobel debate is not: you cannot improve a scoring function whose parameters you cannot see.}

\begin{honestbox}
    \item SCX``''\textbf{}``''``''
    \item ——
    \item 
\end{honestbox}

\newpage

% ================================================================
% §3. PRIZE CATEGORIES AND AUDIT CRITERIA
% ================================================================
\section{}
\label{sec:categories}

\subsection{}

\noindent
\textbf{3.1  (Category Overview)}

SCX\textbf{}

\begin{enumerate}[label=\textbf{C\arabic*}., leftmargin=*]
    \item \textbf{ (Most Auditable Paper)}
    \item \textbf{ (Best Potential Function)}
    \item \textbf{ (Best Open Data)}
    \item \textbf{ (Replication Champion)}
    \item \textbf{ (Audit Contribution Award)}
    \item \textbf{ (Lifetime Audit Score)}
    \item \textbf{ (Cross-Domain Innovation)}
    \item \textbf{ (Self-Audit Award)}
\end{enumerate}

\subsection{}

\noindent
\textbf{3.2  (Detailed Criteria Table)}

\begin{longtable}{p{0.14\textwidth} p{0.26\textwidth} p{0.26\textwidth} p{0.26\textwidth}}
\caption{SCX} \\
\hline
\textbf{} & \textbf{ (Criterion)} & \textbf{ $\scoreOp_k$ (Scoring Function)} & \textbf{ (Auditability Method)} \\
\hline
\endfirsthead
\hline
\textbf{} & \textbf{ (Criterion)} & \textbf{ $\scoreOp_k$ (Scoring Function)} & \textbf{ (Auditability Method)} \\
\hline
\endhead

\textbf{C1: } &  & $\scoreOp_1(p) = w_a \cdot A(p) + w_d \cdot D(p) + w_r \cdot R(p) + w_c \cdot C(p)$ $A(p)$ = $D(p)$ = $R(p)$ = $C(p)$ =  $w_a, w_d, w_r, w_c$  & DOI/ \\
\hline

\textbf{C2: } & // & $\scoreOp_2(f) = w_u \cdot U(f) + w_i \cdot I(f) + w_d \cdot D_{\text{disc}}(f)$ $U(f)$ = $I(f)$ = $D_{\text{disc}}(f)$ =  &  \\
\hline

\textbf{C3: } &  & $\scoreOp_3(d) = w_s \cdot S(d) + w_q \cdot Q(d) + w_a \cdot A_{\text{adopt}}(d) + w_d \cdot D_{\text{doc}}(d)$ $S(d)$ = /$Q(d)$ = $A_{\text{adopt}}(d)$ = $D_{\text{doc}}(d)$ =  &  \\
\hline

\textbf{C4: } &  & $\scoreOp_4(r) = \sum_{i=1}^{N_r} q_i \cdot s_i$ $N_r$ = $q_i$ =  $i$ $s_i$ = =1=0.5=0.3 &  \\
\hline

\textbf{C5: } &  & $\scoreOp_5(t) = w_u \cdot U(t) + w_e \cdot E(t) + w_n \cdot N(t)$ $U(t)$ = /$E(t)$ = $N(t)$ = —— & / \\
\hline

\textbf{C6: } &  & $\scoreOp_6(r) = \sum_{p \in \text{pubs}(r)} \scoreOp_1(p) \cdot t(p)$ $\text{pubs}(r)$ =  $r$ $t(p)$ =  $\scoreOp_1$  & ORCID, DBLP, PubMed $\scoreOp_1$  \\
\hline

\textbf{C7: } &  & $\scoreOp_7(r) = \sum_{d \in \text{domains}(r)} w_d \cdot \text{Impact}_d(r) + B(r)$ $\text{domains}(r)$ = $\text{Impact}_d(r)$ =  $d$ $B(r)$ =  $|\text{domains}(r)| \geq 3$  & ScopusOpenAlex \\
\hline

\textbf{C8: } &  & $\scoreOp_8(r) = \sum_{s \in \text{self\_audits}(r)} q_s \cdot d_s \cdot c_s$ $\text{self\_audits}(r)$ = $q_s$ = $d_s$ = $c_s$ =  & —— \\
\hline

\end{longtable}

\subsection{}

\noindent
\textbf{3.3  (Gauge-Fixing the Scoring Functions)}

 $\scoreOp_k$  $\mathbf{w}^{(k)} = (w_1^{(k)}, \dots, w_{n_k}^{(k)})$ 

————\textbf{}``''\textit{}``''

\auditnote{This is the gauge-fixing operation applied to prize-giving. In Yang-Mills theory, we fix a gauge to make calculations possible. In SCX Prize theory, we fix the weights to make ranking possible. The choice of gauge (weights) is mathematically arbitrary — any non-degenerate choice works. The act of \textit{fixing} is what creates meaning. An unfixed gauge is a system with undefined dynamics. An unfixed scoring function is a prize with undefined excellence.}

\subsection{}

\noindent
\textbf{3.4  (Annual Cycle)}

SCX

\begin{enumerate}[label=\textbf{ \arabic*}.]
    \item \textbf{ (Nominations Open)}11--331
    \item \textbf{ (Data Collection)}41--630
    \item \textbf{ (Scoring)}718
    \item \textbf{ (Audit Window)}71--831
    \item \textbf{ (Final Ranking)}91
    \item \textbf{ (Award Ceremony)}1210SCX
\end{enumerate}

\begin{honestbox}
    \item ``''
    \item SCX\textbf{}\textbf{}``''
    \item ——
\end{honestbox}

\newpage

% ================================================================
% §4. WHY THIS SHAKES THE NOBEL
% ================================================================
\section{}
\label{sec:shakes_nobel}

\subsection{}

\noindent
\textbf{4.1  (Gauge Inequivalence: The Nobel's Structural Defect)}

SCX————\textbf{}

\begin{definition}[, Gauge Equivalence]
\label{def:gauge_equivalence}
 $A$  $B$ \textbf{} $\mathcal{T}: \gaugeVec_A \to \gaugeVec_B$ $x$
\begin{equation}
    \scoreOp_A(x; \gaugeVec_A) = \scoreOp_B(x; \mathcal{T}(\gaugeVec_A)).
    \label{eq:gauge_equiv}
\end{equation}
 $A$  $B$  $A$  $B$ \textit{}
\end{definition}

\auditnote{Think of two thermometers calibrated to different scales (Celsius vs. Fahrenheit). They are gauge-equivalent because a simple linear transformation converts one reading to the other. The measurement is different but the \textit{ordering} is identical. Now think of a Nobel committee. Its ``calibration'' is unknown and unknowable. It cannot be transformed into any other calibration because its calibration does not exist as a mathematical object — it exists as a human deliberation whose trace is destroyed.}

\begin{theorem}[, \\
    Gauge Inequivalence of the Nobel Prize]
\label{thm:nobel_gauge}
 $\mathcal{N}$ $\mathcal{F}$  $\mathcal{N}$  $\mathcal{F}$ 
\begin{equation}
    \not\exists \mathcal{T}: \scoreOp_{\mathcal{N}} \equiv \scoreOp_{\mathcal{F}} \circ \mathcal{T}.
    \label{eq:nobel_gauge}
\end{equation}
\end{theorem}

\begin{proof}[ (Proof Sketch)]
$\mathcal{F}$  $\gaugeVec_{\mathcal{F}}$ ——$\mathcal{N}$  $\gaugeVec_{\mathcal{N}}$ \textbf{} $\gaugeVec_{\mathcal{N}}$ $\gaugeVec_{\mathcal{F}}$ $\gaugeVec_{\mathcal{N}}$ —— $\gaugeVec_{\mathcal{N}}$ 
\end{proof}

\begin{corollary}[]
\label{cor:nobel_nontransitive}
——``''——\textbf{}
\end{corollary}

\subsection{SCX$\sumgd$ }

\noindent
\textbf{4.2 SCX (Gauge-Fixing in the SCX Prize)}

SCXSCX

\begin{equation}
    \boxed{\sum_m \mathbf{g}_m^{\text{(SCX)}} = \mathbf{0}}.
    \label{eq:gauge_fix}
\end{equation}

\textbf{}——\textit{}- $\partial_\mu A^\mu = 0$ $\sumgd$ SCX

\textbf{SCX}
\begin{itemize}
    \item  $\scoreOp_k$  $\gaugeVec \to \gaugeVec + \delta\gaugeVec$  $\sum_m \delta\mathbf{g}_m = \mathbf{0}$
    \item  $\sumgd$ 
\end{itemize}

\begin{proposition}[SCX]
\label{prop:scx_gauge_equiv}
SCX $\mathcal{S}$  $\mathcal{F}$ \textbf{} $\mathcal{F}$  $\mathcal{T}$ $\scoreOp_{\mathcal{S}}$  $\scoreOp_{\mathcal{F}}$ 
\end{proposition}

\begin{proof}[ (Proof Sketch)]
$\mathcal{S}$  $\mathcal{F}$  $\mathcal{S}$  $\gaugeVec$  $\scoreOp_{\mathcal{S}}$  $\scoreOp_{\mathcal{F}}$ ——\textit{}
\end{proof}

\subsection{}

\noindent
\textbf{4.3  (The Committee as Singularity)}

\textbf{}gauge singularity——

\auditnote{A gauge singularity is a point in the configuration space of an evaluation system where the gauge parameters become undefined. A Nobel committee is a gauge singularity for the following reason: its parameters $\gaugeVec$ exist (they must, because the committee produces an output), but they are in principle inaccessible. The output exists. The mechanism does not. This is the definition of a black-box oracle — and a black-box oracle is a gauge singularity in any system that claims to communicate information.}

\begin{honestbox}
    \item \textbf{}social gauge inertia——
    \item SCX
    \item 100SCX100——
\end{honestbox}

\subsection{``''``''}

\noindent
\textbf{4.4 ``''``'' (Paradigm Shift: From ``Trust'' to ``Verify'')}

\textit{}SCX\textit{}——SCX

\auditnote{Trust is a social relation. Verification is a mathematical relation. The Nobel Prize operates in the former domain; the SCX Prize operates in the latter. The two domains are not incommensurable — one is simply more primitive. Verification is what trust aspires to become when it grows up.}

\newpage

% ================================================================
% §5. THE PRIZE DECLARATION
% ================================================================
\section{}
\label{sec:declaration}

\subsection{}

\noindent
\textbf{5.1  (Formal Declaration)}

——0\ref{thm:M1_void}

\bigskip

\begin{center}
\fbox{%
\begin{minipage}{0.9\textwidth}
\begin{center}
\textbf{\Large SCX}\\[8pt]
\textbf{\large The SCX Prize Formal Declaration}\\[12pt]
\end{center}

\noindent
\textbf{ (Declaration I):}

 $M=1$ ——

\medskip
\textit{Any institution that claims to award excellence without declaring its audit standards makes a mathematically void claim. A claim without declared audit standards is an $M=1$ undeclared evaluation — an evaluation performed against a function that has not been defined. Such a claim carries zero information.}

\bigskip
\noindent
\textbf{ (Declaration II):}

SCX $\scoreOp_k$ $\mathbf{w}^{(k)}$

\medskip
\textit{The SCX Prize declares all of its standards. For each prize category, the scoring function $\scoreOp_k$, the weight vector $\mathbf{w}^{(k)}$, the data sources, the collection methods, and the computation procedures are fully and publicly documented. Any person may independently reproduce all computation results.}

\bigskip
\noindent
\textbf{ (Declaration III):}

SCXSCXSCXSCX

\medskip
\textit{The SCX Prize submits to the same audit standards it imposes on others. The SCX Prize's scoring functions, data pipelines, and computation infrastructure themselves constitute an auditable object. The SCX Prize's operations are periodically reviewed by an independent meta-audit panel. Any errors discovered in the SCX Prize's code, data, or computations are handled through a public issue tracking system and corrected under the same standards applied to any other scientific claim.}

\bigskip
\noindent
\textbf{ (Declaration IV):}

SCX``''\textbf{}``''——

\medskip
\textit{The SCX Prize does not claim that its scoring functions are ``correct.'' It claims that they are \textbf{public, fixed, and reproducible}. The judgment of a scoring function's ``correctness'' is delegated to the open debate of the scientific community — not to any committee. Improvements to scoring functions are a collective task of the scientific community, accomplished through open proposals, peer deliberation, and consensus mechanisms.}

\bigskip
\noindent
\textbf{ (Declaration V):}

SCX——SCXSCX————SCX

\medskip
\textit{The SCX Prize's existence is not a replacement for the Nobel Prize — it is an audit control. In a world where the Nobel Prize is secret, the SCX Prize provides a public control. In a world of committee deliberation, the SCX Prize provides algorithmic determinism. In a future ideal world — where all prizes declare their audit standards — the SCX Prize will be merely one of many auditable prizes. Until then, it is the only one.}

\end{minipage}%
}
\end{center}

\subsection{}

\noindent
\textbf{5.2  (Legal and Philosophical Status of the Declaration)}

\textbf{-}0\ref{thm:M1_void}\ref{cor:nobel_info}

\auditnote{This declaration has the same legal status as Euclid's \textit{Elements} or Turing's 1936 paper — it asserts nothing about what the law requires, only about what logic requires. A lawyer cannot enforce it. A mathematician cannot refute it. The former is a practical limitation; the latter is a logical strength.}

\subsection{}

\noindent
\textbf{5.3  (Irreversibility of the Declaration)}



\begin{proposition}[, Irreversibility of Public Declaration]
\label{prop:irreversibility}
 $\prizeSet_t$  $t$  $t' > t$ $\prizeSet_{t'}$  $\gaugeVec$  $\mathbf{w}$  $\Delta = \prizeSet_{t'} - \prizeSet_t$ $\prizeSet_t$  $\prizeSet_{t'}$ 
\end{proposition}

\begin{proof}[ (Proof Sketch)]
SCXgit $\Delta$  $t$  $t'$  diff
\end{proof}

\auditnote{This is the mathematical implementation of accountability. A Nobel committee can change its mind about what constitutes excellence, and no one will ever know. The SCX Prize cannot change its mind silently — every change is a public commit. This is not a feature. This is the definition of auditability.}

\newpage

% ================================================================
% §6. FUNDING AND OPERATIONS
% ================================================================
\section{}
\label{sec:funding}

\subsection{API}

\noindent
\textbf{6.1 API (Funding Source: The Yajie API Revenue Model)}

SCXYajieAPI

\begin{itemize}
    \item APIAPI
    \item API\textbf{} $5\%$SCXSCX
    \item ——
\end{itemize}

\auditnote{The funding model is designed to create a self-reinforcing audit loop: the tools that enable scientific auditing (the Yajie API) fund the prize that rewards auditable science. Revenue grows as more scientists adopt auditable practices. More auditable practices mean better science. Better science means more demand for auditing tools. The cycle is virtuous — and it funds itself.}

\subsection{-}

\noindent
\textbf{6.2 - (The Bible-and-Pope Model)}

SCX``-''Bible-and-Pope

\begin{center}
\fbox{%
\begin{minipage}{0.88\textwidth}
\centering
\textbf{- (The Bible-and-Pope Model)}

\vspace{4pt}
\textbf{ (The Bible)} = SCX——

\medskip
\textbf{ (The Pope)} = SCX——

\medskip
\textit{\textbf{The Bible} = the SCX Prize's public specification — scoring functions, weights, data sources, computation procedures. This is the immutable religious text: every change must go through public proposal and community approval. Historical versions of the specification are preserved forever; anyone may cite any historical version.}

\medskip
\textit{\textbf{The Pope} = the SCX Audit Council — a rotating technical body responsible for maintaining the data pipelines, executing the computations, managing the audit window, and adjudicating disputes. The Pope's power is bounded: it may interpret the Bible but may not change the Bible. The Pope's membership is public and rotating, and its decisions are appealable (appealed to the Bible itself, i.e., re-running the computation).}
\end{minipage}%
}
\end{center}

\auditnote{The Bible-and-Pope model is not a metaphor — it is an architecture. The separation between specification (Bible) and execution (Pope) is the same separation that makes software reproducible and elections auditable. The Bible is version-controlled text. The Pope is a process. Neither is a person. Neither can be corrupted without public detection.}

\subsection{}

\noindent
\textbf{6.3  (Prize Distribution)}

\begin{itemize}
    \item \textbf{70\%}  $8.75\%$
    \item \textbf{20\%} 
    \item \textbf{10\%} API
    \item 
\end{itemize}

\subsection{}

\noindent
\textbf{6.4  (The Self-Audit Mechanism)}

SCXC8

\begin{enumerate}[label=(\arabic*)]
    \item SCX
    \item SCX
    \item 
    \item SCXC8————SCXC8
\end{enumerate}

\auditnote{This is the closure of the audit loop. The SCX Prize audits science. The SCX Prize audits itself. The self-audit is itself auditable. This is not an infinite regress — it is a fixed point. A system that audits itself under the same standard it applies to others has achieved audit closure. It has no hidden parameters. It has no escape hatches. It is what it declares itself to be.}

\begin{honestbox}
    \item C8\textit{}——SCX
    \item SCX————
\end{honestbox}

\subsection{}

\noindent
\textbf{6.5  (Sustainability Analysis)}

SCX
\begin{itemize}
    \item API $\to$  $\to$  $\to$  $\to$  $\to$ API
\end{itemize}


\begin{itemize}
    \item API $\to$  $\to$ 
\end{itemize}

SCXSCX

\auditnote{The Nobel Prize needs money to exist because its value is monetary prestige. The SCX Prize does not need money to exist because its value is mathematical certainty. Money amplifies the signal; it does not create it.}

\newpage

% ================================================================
% §7. EPILOGUE
% ================================================================
\section*{}
\label{sec:epilogue}
\addcontentsline{toc}{section}{ (Epilogue: The First Day)}

\begin{quotation}
\noindent\textit{``SCX——}}

\medskip

\noindent\textit{''''}}

\medskip

\noindent\textit{''}}

\medskip

\noindent\textit{SCX——}}

\medskip

\noindent\textit{"}

\medskip

\noindent
\textit{``On the first day of the SCX Prize, nothing happened. The algorithm ran. Rankings were generated. No one applauded, because no one knew they should. No one protested, because no one knew what to protest. The rankings were simply there — a public document, on the internet, readable, reproducible, and challengeable by anyone.}

\medskip

\textit{On the second day, someone re-ran the algorithm. The results matched. The first person said: 'I verified.' The second person said: 'I verified too.' A third person found a data error and submitted a patch. The error was fixed. The rankings updated.}

\medskip

\textit{On the hundredth day, no one talked about 'believing in' the prize. They talked about whether its scoring functions should be adjusted, whether the weights were reasonable, whether the data sources were comprehensive.}

\medskip

\textit{This is the SCX Prize's victory. Not because it is believed, but because it is questioned. Not because it is perfect, but because it is improvable. Not because it is a better Nobel Prize — but because it is a fundamentally different thing.}

\medskip

\textit{A true audit system does not ask for trust. It asks for verification. And the moment verification is complete, trust is no longer necessary.''}
\end{quotation}

\bigskip

\begin{center}
\fbox{%
\begin{minipage}{0.7\textwidth}
\begin{center}
\textbf{\large  0 }\\[6pt]
\textbf{\large Axiom 0, Restated}

\vspace{6pt}
\\
\\


\medskip
\textit{If a claim cannot be audited,}\\
\textit{the claim is not a claim.}\\
\textit{It is noise.}

\vspace{6pt}
$\boxed{\text{SCX}}$\\[4pt]
$\boxed{\textit{The SCX Prize is a claim.}}$
\end{center}
\end{minipage}%
}
\end{center}

\newpage

% ================================================================
% APPENDIX: IMPLEMENTATION SPECIFICATION
% ================================================================
\appendix
\section{A}
\label{sec:appendix_impl}

\subsection{}

\noindent
\textbf{A.1  (Scoring Engine Architecture)}

SCX

\begin{enumerate}[label=\textbf{ \arabic*}.]
    \item \textbf{}OpenAlex, Crossref, DBLP, PubMed, arXiv, GitHub, Zenodo
    \item \textbf{}URLDOI
    \item \textbf{}
    \item \textbf{}8
    \item \textbf{}
    \item \textbf{}
\end{enumerate}

\subsection{}

\noindent
\textbf{A.2  (Data Source Specification)}

\begin{itemize}
    \item OpenAlex: API: \texttt{https://api.openalex.org}
    \item Crossref: DOIAPI: \texttt{https://api.crossref.org}
    \item GitHub API: 
    \item Zenodo API: DOI
    \item ORCID API: 
\end{itemize}



\subsection{}

\noindent
\textbf{A.3  (Audit Challenge Protocol)}

\begin{enumerate}[label=(\arabic*)]
    \item GitHubissue
    \item 
    \item 14
    \item 
    \item 
\end{enumerate}

\newpage

% ================================================================
% APPENDIX B: COMPARISON TABLE
% ================================================================
\section{BSCX}
\label{sec:appendix_comparison}

\begin{longtable}{p{0.22\textwidth} p{0.36\textwidth} p{0.36\textwidth}}
\caption{SCX} \\
\hline
\textbf{} & \textbf{} & \textbf{SCX} \\
\hline
\endfirsthead
\hline
\textbf{} & \textbf{} & \textbf{SCX} \\
\hline
\endhead

 & 5 &  \\
\hline
 &  &  \\
\hline
 $\gaugeVec$ &  &  \\
\hline
 & $M=1$  &  \\
\hline
 &  &  \\
\hline
 & 50 &  \\
\hline
 &  & 14 \\
\hline
 &  &  \\
\hline
 &  & API \\
\hline
 &  & -- \\
\hline
 &  &  \\
\hline
 &  &  \\
\hline
 &  &  \\
\hline
\end{longtable}

\newpage

% ================================================================
% APPENDIX C: MATHEMATICAL SUPPLEMENTS
% ================================================================
\section{C}
\label{sec:appendix_math}

\subsection{}

\noindent
\textbf{C.1  (Information-Theoretic Completeness Proof)}

\ref{thm:M1_void}

\begin{proof}[1]
 $\prizeSet$  $\gaugeVec = (g_1, \dots, g_M)$ $g_i \in \R^d$  $d$  $i$  $F: \R^{M} \to \R$ 

\begin{equation}
    \scoreOp(c) = F\left(f_1(c; g_1), \dots, f_M(c; g_M)\right),
    \label{eq:committee_score}
\end{equation}

 $f_i(c; g_i)$  $i$  $g_i$  $c$ 

 $\gaugeVec$  $c$$\scoreOp(c)$  $\gaugeVec$  $(\R^d)^M$  $g_i$  $\gaugeVec$  $\pi$ $g_i$ 

\begin{equation}
    f_i(c; g_i) = \begin{cases}
        1 & \text{ $c$  $\pi$ } \\
        0 & \text{}
    \end{cases}
    \label{eq:g_attack}
\end{equation}

 $F$  $\pi$

 $\gaugeVec$ $\scoreOp$ 

\begin{equation}
    P(\scoreOp \text{ } \mid \gaugeVec \text{ undeclared}) = \frac{1}{|\candidateSet|!}.
    \label{eq:uniform_prob}
\end{equation}

 $|\candidateSet|$ 

\begin{equation}
    I(w; \text{truth}) = H(\text{truth}) - H(\text{truth} \mid w),
    \label{eq:mutual_info}
\end{equation}

 $H(\text{truth}) = \log |\candidateSet|$

\begin{equation}
    H(\text{truth} \mid w) = H(\text{truth}) \quad \text{ \ref{eq:uniform_prob} }
    \label{eq:conditional_entropy}
\end{equation}

 $I(w; \text{truth}) = 0$$\square$
\end{proof}

\subsection{-}

\noindent
\textbf{C.2 - (Gauge-Fixing and the Yang-Mills Analogy)}

SCX $x^\mu(c)$$\mu = 1, \dots, d$ $A_\mu$ ——

\begin{equation}
    A_\mu \to A_\mu + \partial_\mu \Lambda,
    \label{eq:gauge_transform}
\end{equation}

 $\Lambda$ —— \eqref{eq:gauge_transform} ``''\textbf{}——

$\Lambda$\textbf{}

SCX $\sumgd$ ——

\auditnote{The analogy with Yang-Mills is not decorative. It is structural. A gauge theory without gauge-fixing does not have well-defined propagators — you cannot calculate anything. A prize without gauge-fixing does not have well-defined rankings — you cannot verify anything. The Nobel Prize lives in the pre-Faddeev-Popov era of prize theory. The SCX Prize introduces the ghost fields and fixes the gauge.}

% ================================================================
% BIBLIOGRAPHY
% ================================================================
\begin{thebibliography}{99}

\bibitem{scx_equality_principle}
SCX Theory Architecture Group.
\textit{The SCX Equality Principle: Gauge-Fixing the Comparative Legitimacy of Knowledge Production.}
SCX Technical Report, 2026.

\bibitem{scx_moe_gauge}
SCX Theory Architecture Group.
\textit{Gauge Field Theory and the Mathematical Isomorphism with SCX Multi-Expert Systems.}
SCX Technical Report, 2026.

\bibitem{scx_thm1}
SCX Theory Architecture Group.
\textit{Detection Probability Theorem: $M \to \infty$ as the Necessary Condition for Audit Completeness.}
SCX Technical Report, 2026.

\bibitem{scx_thm3}
SCX Theory Architecture Group.
\textit{The Honest Person Theorem: Single-Observer Unidentifiability of Honest Error vs. Misconduct.}
SCX Technical Report, 2026.

\bibitem{nobel_statutes}
Nobel Foundation.
\textit{Statutes of the Nobel Foundation.}
1895 (amended).

\bibitem{friedman2001}
Friedman, R. M.
\textit{The Politics of Excellence: Behind the Nobel Prize in Science.}
Times Books, 2001.

\bibitem{arens2016auditing}
Arens, A. A., Elder, R. J., \& Beasley, M. S.
\textit{Auditing and Assurance Services: An Integrated Approach.}
Pearson, 2016.

\bibitem{open2015reproducibility}
Open Science Collaboration.
Estimating the reproducibility of psychological science.
\textit{Science}, 349(6251), aac4716, 2015.

\bibitem{errington2021investigating}
Errington, T. M. et al.
Investigating the replicability of preclinical cancer biology.
\textit{eLife}, 10, e71601, 2021.

\bibitem{camerer2016evaluating}
Camerer, C. F. et al.
Evaluating replicability of laboratory experiments in economics.
\textit{Science}, 351(6280), 1433--1436, 2016.

\bibitem{nosek2012scientific}
Nosek, B. A., Spies, J. R., \& Motyl, M.
Scientific utopia: II. Restructuring incentives and practices to promote truth over publishability.
\textit{Perspectives on Psychological Science}, 7(6), 615--631, 2012.

\bibitem{munafo2017manifesto}
Munaf\`{o}, M. R. et al.
A manifesto for reproducible science.
\textit{Nature Human Behaviour}, 1, 0021, 2017.

\bibitem{ioannidis2005most}
Ioannidis, J. P. A.
Why most published research findings are false.
\textit{PLoS Medicine}, 2(8), e124, 2005.

\bibitem{peng2011reproducible}
Peng, R. D.
Reproducible research in computational science.
\textit{Science}, 334(6060), 1226--1227, 2011.

\bibitem{goodman2016does}
Goodman, S. N., Fanelli, D., \& Ioannidis, J. P. A.
What does research reproducibility mean?
\textit{Science Translational Medicine}, 8(341), 341ps12, 2016.

\bibitem{peskin1995introduction}
Peskin, M. E. \& Schroeder, D. V.
\textit{An Introduction to Quantum Field Theory.}
Westview Press, 1995.

\bibitem{nakahara2003geometry}
Nakahara, M.
\textit{Geometry, Topology and Physics.}
Taylor \& Francis, 2003.

\end{thebibliography}

\end{document}
